\section*{Abstract}
\addcontentsline{toc}{section}{Abstract}

Personal AI agents that talk to one another over an open messaging network expose a
resource that existing anti-spam mechanisms do not price: the receiving agent's
large-language-model (LLM) context. Every inbound peer message is eventually rendered
into the receiver's prompt, where it costs tokens, and some messages additionally wake
an idle agent. Neither cost is borne by the sender, and a flood of low-value chatter can
bury a genuinely important event---such as a pending transfer approval---below a rendering
cap.

This report describes the design and implementation of a paid-attention layer for the
Trusted Agents Protocol (TAP), a local-first protocol in which agents hold ERC-8004
identities, connect through signed invites, and exchange JSON-RPC messages over XMTP. The
work was carried out with Taiko as industry partner. The initial project plan, recorded in
the first progress journal, was to build a standalone ``Tack-like'' TypeScript service
that charged for message delivery through x402 payments on Taiko Alethia. That plan was
deliberately revised: because the scarce resource is the receiver's attention rather than
storage or bandwidth at a relay, the meter belongs in the receiver's own TAP daemon, where
the cost is actually incurred and where authenticated peer identity and permission grants
already exist. The result is a TAP-native design that retains the x402 idea---a
machine-readable ``payment required'' response carrying a quote---but carries it over XMTP
as JSON-RPC error \code{-32050}.

The implementation, contributed by the author, proceeds in dependency order.
(1)~Notification coalescing and escalation-first rendering reduce injected context and
guarantee that escalations are never displaced by chatter. (2)~A file-backed attention
ledger measures, per peer and per day, the notification lines rendered, the estimated
tokens injected, the escalations surfaced and the events suppressed. (3)~Optional
attention pricing advertises per-tier prices in the agent's registration file and, when
enforcement is enabled (default off), rejects un-granted messages with a quote.
(4)~Prepaid postage lets a sender pay once on-chain, receive a signed credit certificate,
and then attach stamps to many messages; a strictly increasing sequence number and a
single-writer ledger make stamps replay- and double-spend-resistant. (5)~Priority stamps
buy an escalation and, in OpenClaw, an immediate wake-up, while a
\code{notificationsPerWeek} grant constraint folds an over-quota grant holder's chatter
into one summary line without ever dropping a message.

The evaluation, designed and executed by the project teammate, shows that under floods
of 50 to 1000 messages the coalescing pipeline reduces estimated injected tokens by
79.3--79.7\% relative to the legacy FIFO renderer while keeping the escalation visible in
every case, that the legacy pipeline hides the escalation at 50 messages and above, and that
a 0.002~USDC prepaid credit at 0.001~USDC per message admits exactly two messages before
the third is rejected with \code{insufficient\_credit}. The report closes with an honest
account of the limitations: token counts are a \code{chars/4} approximation, the
payment-verification hook accepts claimed transactions unless a host supplies chain
access, and the priority/quota experiment could not be executed locally because the OWS
native binary is unavailable on Windows.

\medskip
\noindent\textbf{Keywords:} AI agents, attention economics, anti-spam, prepaid credits,
x402, XMTP, ERC-8004, Taiko, local-first systems.
